\paragraph{An accuracy-independent gradient prefix}
\label{sec:gradient-prefix}

Fix any finite prefix of $K$ accepted steps, with the endpoint convention
for $\tau_K$ specified above. All estimates below concern this finite
prefix; termination has not yet been assumed. Set
\[
 G_*=(\Delta^2Q)^{1/3}>0,\qquad c_r=\frac{51}{3^{2/3}}.
\]
For a boundary step, \eqref{eq:boundary} and
\eqref{eq:descent-budget} give
\begin{align}
 \norm{\nabla f(x_{\tau_{j+1}})}
 &\le51\sigma_{\tau_j}r_{\tau_j}^2\nonumber\\
 &=51\sigma_{\tau_j}^{1/3}
             (\sigma_{\tau_j}r_{\tau_j}^3)^{2/3}\nonumber\\
 &\le c_r\sigma_{\tau_j}^{1/3}
       [f(x_{\tau_j})-f(x_{\tau_{j+1}})]^{2/3}.
 \label{eq:boundary-reset}
\end{align}
Every function decrease in this expression is at most $\Delta$.
The global scale bound \eqref{eq:scale-global} therefore shows that
each boundary step produces a gradient norm no larger than
$c_r(C_EQ_0)^{1/3}\Delta^{2/3}$. Interior steps multiply the
current gradient norm by at most $\rho=11/16$. Induction on $j$,
starting from $x_{\tau_0}=x_0$, consequently gives
\begin{equation}\label{eq:global-gradient-bound}
 \norm{\nabla f(x_{\tau_j})}\le H_0
 :=\max\left\{\norm{\nabla f(x_0)},
                c_r(C_EQ_0)^{1/3}\Delta^{2/3}\right\},
 \quad 0\le j\le K.
\end{equation}

For $j\ge j_\sigma$, the sharper bound
$\sigma_{\tau_j}\le C_PQ$ in Proposition~\ref{prop:potential}
and \eqref{eq:boundary-reset} show that a boundary step produces
a gradient norm no larger than
\[
 H:=c_rC_P^{1/3}G_*.
\]
Define the nonnegative integers
\begin{equation}\label{eq:gradient-prefix-length}
 \ell:=\left\lceil\frac{\lp(H_0/H)}{\log(16/11)}\right\rceil,
 \qquad s:=j_\sigma+\ell.
\end{equation}
If $K<s$, this prefix already has fewer than $s$ accepted steps.
Suppose instead that $K\ge s$. If all $\ell$ steps with accepted
indices $j_\sigma,\ldots,s-1$ are interior, repeated use of
\eqref{eq:interior} gives
\[
 \norm{\nabla f(x_{\tau_s})}
 \le\rho^\ell\norm{\nabla f(x_{\tau_{j_\sigma}})}
 \le\rho^\ell H_0\le H.
\]
The last inequality follows directly from the ceiling in
\eqref{eq:gradient-prefix-length}; it also holds when $\ell=0$,
since then $H_0\le H$.
If one of these steps is a boundary step, its next-point gradient
norm is at most $H$. Every later interior step contracts this norm,
and every later boundary step again bounds it by $H$. Thus in both
cases
\begin{equation}\label{eq:post-prefix-gradient}
 \norm{\nabla f(x_{\tau_s})}\le H.
\end{equation}

We next bound the prefix length explicitly in terms of the analysis
parameters. Since $c_rC_P^{1/3}>1$ and $C_E<C_P$, the definitions imply
\[
 \frac{H_0}{H}
 \le\max\left\{
       \frac{\norm{\nabla f(x_0)}}{G_*},
       \left(\frac{Q_0}{Q}\right)^{1/3}\right\}.
\]
Using $\lceil z\rceil\le z+1$ for $z\ge0$ and
$\log_+\max\{a,b\}\le\log_+a+\log_+b$ for $a\ge0$, $b>0$,
we obtain
\[
 \ell\le1+\frac{1}{\log(16/11)}
 \left(\lp\frac{\norm{\nabla f(x_0)}}{G_*}
                   +\frac13\lp\frac{Q_0}{Q}\right).
\]
Together with the bound for $j_\sigma$ from
Proposition~\ref{prop:potential}, this proves
\begin{equation}\label{eq:prefix-bound}
 s=O\left(1+\lp\frac{Q_0}{Q}
                 +\lp\frac{\norm{\nabla f(x_0)}}{G_*}\right).
\end{equation}
In particular, $s$ does not depend on $\epsilon$ or $\delta$.
These indices are used only in the proof; the algorithm does not
compute them or restart at $s$.

\paragraph{Sums of gradient norms and negative curvature}
\label{sec:stationarity-sums}

Suppose $K\ge s+m$ for some integer $m\ge1$, and let
\[
 \mathcal J=\{s,s+1,\ldots,s+m-1\},\qquad
 T=\sum_{j\in\mathcal J}\sigma_{\tau_j},\qquad
 S=mL_0+L_1^{3/2}\Delta.
\]
Partition $\mathcal J$ into the boundary indices
$\mathcal J_\partial:=\{j\in\mathcal J:r_{\tau_j}=R_{\tau_j}\}$
and the interior indices $\mathcal J_\circ:=\mathcal J\setminus
\mathcal J_\partial$. The scale estimate \eqref{eq:suffix-scale}
and the descent budget \eqref{eq:descent-budget} imply
\begin{equation}\label{eq:suffix-budgets}
 \begin{aligned}
 T&\le C_TS,\\
 \sum_{j\in\mathcal J}\sigma_{\tau_j}r_{\tau_j}^3
 &\le\frac13\sum_{j\in\mathcal J}
          [f(x_{\tau_j})-f(x_{\tau_{j+1}})]
 \le\frac{\Delta}{3}.
 \end{aligned}
\end{equation}
Also $S\ge Q$ because $m\ge1$.

Both cases in \eqref{eq:boundary}--\eqref{eq:interior} imply
the single recurrence
\[
 \norm{\nabla f(x_{\tau_{j+1}})}
 \le\rho\norm{\nabla f(x_{\tau_j})}
       +51\mathbf1_{\{j\in\mathcal J_\partial\}}
            \sigma_{\tau_j}r_{\tau_j}^2,
 \qquad j\in\mathcal J.
\]
Indeed, on an interior step the second term vanishes, while on a
boundary step \eqref{eq:boundary} bounds the left-hand side by
the second term alone. Summing this recurrence and retaining the
two endpoint terms gives
\[\begin{aligned}
 &(1-\rho)\sum_{j\in\mathcal J}\norm{\nabla f(x_{\tau_j})}\\
 &\quad\le\norm{\nabla f(x_{\tau_s})}
          -\norm{\nabla f(x_{\tau_{s+m}})}
          +51\sum_{j\in\mathcal J_\partial}
                 \sigma_{\tau_j}r_{\tau_j}^2\\
 &\quad\le H+51\sum_{j\in\mathcal J_\partial}
                 \sigma_{\tau_j}r_{\tau_j}^2.
\end{aligned}\]
The last line uses \eqref{eq:post-prefix-gradient} and nonnegativity
of the terminal gradient norm.
For the remaining sum, H\"older's inequality with exponents $3/2$
and $3$ yields
\begin{align*}
 \sum_{j\in\mathcal J_\partial}\sigma_{\tau_j}r_{\tau_j}^2
 &=\sum_{j\in\mathcal J_\partial}
       (\sigma_{\tau_j}r_{\tau_j}^3)^{2/3}
       \sigma_{\tau_j}^{1/3}\\
 &\le\left(\sum_{j\in\mathcal J_\partial}
                    \sigma_{\tau_j}r_{\tau_j}^3\right)^{2/3}
       \left(\sum_{j\in\mathcal J_\partial}
                    \sigma_{\tau_j}\right)^{1/3}\\
 &\le(\Delta/3)^{2/3}T^{1/3}
 \le3^{-2/3}C_T^{1/3}\Delta^{2/3}S^{1/3}.
\end{align*}
Empty boundary sums equal zero, so this argument also covers a
suffix consisting entirely of interior steps. Since
$H=c_rC_P^{1/3}\Delta^{2/3}Q^{1/3}$ and $Q\le S$, we conclude that
\begin{equation}\label{eq:sum-g}
 \begin{aligned}
 \sum_{j\in\mathcal J}\norm{\nabla f(x_{\tau_j})}
 &\le C_g\Delta^{2/3}S^{1/3},\\
 C_g&:=\frac{c_rC_P^{1/3}+51\cdot3^{-2/3}C_T^{1/3}}{1-\rho}.
 \end{aligned}
\end{equation}

For the curvature sum, first use the interior estimate in
\eqref{eq:interior}, followed by Cauchy--Schwarz:
\begin{align*}
 &\sum_{j\in\mathcal J_\circ}
       [-\lambda_{\min}(\nabla^2f(x_{\tau_j}))]_+\\
 &\quad\le2\sum_{j\in\mathcal J_\circ}
                 \sqrt{\sigma_{\tau_j}\norm{\nabla f(x_{\tau_j})}}\\
 &\quad\le2\left(\sum_{j\in\mathcal J_\circ}\sigma_{\tau_j}\right)^{1/2}
             \left(\sum_{j\in\mathcal J_\circ}
                          \norm{\nabla f(x_{\tau_j})}\right)^{1/2}\\
 &\quad\le2\sqrt{C_TC_g}\,\Delta^{1/3}S^{2/3}.
\end{align*}
For boundary indices, \eqref{eq:boundary} and H\"older's inequality
with exponents $3$ and $3/2$ instead give
\begin{align*}
 &\sum_{j\in\mathcal J_\partial}
       [-\lambda_{\min}(\nabla^2f(x_{\tau_j}))]_+\\
 &\quad\le32\sum_{j\in\mathcal J_\partial}\sigma_{\tau_j}r_{\tau_j}\\
 &\quad=32\sum_{j\in\mathcal J_\partial}
       (\sigma_{\tau_j}r_{\tau_j}^3)^{1/3}\sigma_{\tau_j}^{2/3}\\
 &\quad\le32\left(\sum_{j\in\mathcal J_\partial}
                     \sigma_{\tau_j}r_{\tau_j}^3\right)^{1/3}
              \left(\sum_{j\in\mathcal J_\partial}
                     \sigma_{\tau_j}\right)^{2/3}\\
 &\quad\le32\cdot3^{-1/3}C_T^{2/3}\Delta^{1/3}S^{2/3}.
\end{align*}
Adding the two estimates proves
\begin{equation}\label{eq:sum-mu}
 \sum_{j\in\mathcal J}[-\lambda_{\min}(\nabla^2f(x_{\tau_j}))]_+
 \le C_{\rm curv}\Delta^{1/3}S^{2/3},
\end{equation}
where the universal numerical constant is
\[
 C_{\rm curv}:=2\sqrt{C_TC_g}+32\cdot3^{-1/3}C_T^{2/3}.
\]
In particular, the curvature estimate concerns the current accepted
points. It does not require a Taylor estimate of the next-point
Hessian at a trial scale that may underestimate the safe scale.

\paragraph{The four-term bound and finite termination}
\label{sec:termination}

At every accepted trial the stopping test has failed. Therefore,
for each $j\in\mathcal J$, at least one of
$\norm{\nabla f(x_{\tau_j})}>\epsilon$ and
$[-\lambda_{\min}(\nabla^2f(x_{\tau_j}))]_+>\delta$ holds, and
\[
 1<\frac{\norm{\nabla f(x_{\tau_j})}}{\epsilon}
     +\frac{[-\lambda_{\min}(\nabla^2f(x_{\tau_j}))]_+}{\delta}.
\]
Sum over $\mathcal J$, divide by $m$, and apply
\eqref{eq:sum-g} and \eqref{eq:sum-mu}. This gives
\[
 1<\frac{C_g\Delta^{2/3}S^{1/3}}{\epsilon m}
       +\frac{C_{\rm curv}\Delta^{1/3}S^{2/3}}{\delta m}.
\]
For nonnegative $a,b$ and $0<p<1$,
$(a+b)^p\le a^p+b^p$: when $a+b>0$, divide by $(a+b)^p$
and use $t^p\ge t$ on $0\le t\le1$ for
$t=a/(a+b)$ and $1-t=b/(a+b)$; the case $a=b=0$ is immediate.
Apply this inequality with $p=1/3$ and $p=2/3$ to
$S=mL_0+L_1^{3/2}\Delta$. The two resulting expressions are
\begin{align*}
 \frac{\Delta^{2/3}S^{1/3}}{\epsilon m}
 &\le\frac{\Delta^{2/3}L_0^{1/3}}{\epsilon m^{2/3}}
       +\frac{\sqrt{L_1}\Delta}{\epsilon m},\\
 \frac{\Delta^{1/3}S^{2/3}}{\delta m}
 &\le\frac{\Delta^{1/3}L_0^{2/3}}{\delta m^{1/3}}
       +\frac{L_1\Delta}{\delta m}.
\end{align*}
With $C_*:=\max\{1,C_g,C_{\rm curv}\}$ we have therefore proved
\begin{equation}\label{eq:average}
 1<C_*\left[
 \frac{\Delta^{2/3}L_0^{1/3}}{\epsilon m^{2/3}}
 +\frac{\sqrt{L_1}\Delta}{\epsilon m}
 +\frac{\Delta^{1/3}L_0^{2/3}}{\delta m^{1/3}}
 +\frac{L_1\Delta}{\delta m}\right].
\end{equation}

To make the contradiction explicit, set $q=(8C_*)^3$ and define
\begin{equation}\label{eq:tail-length}
 m_*:=\left\lceil q\left[1+\Delta\left(
       \frac{\sqrt{L_0}}{\epsilon^{3/2}}
       +\frac{\sqrt{L_1}}{\epsilon}
       +\frac{L_0^2}{\delta^3}
       +\frac{L_1}{\delta}\right)\right]\right\rceil.
\end{equation}
If $m\ge m_*$, each of the four ratios of a complexity term in
the brackets of \eqref{eq:tail-length} to $m$ is at most $q^{-1}$.
The four terms in brackets in \eqref{eq:average} are respectively
the $2/3$ power, the first power, the $1/3$ power, and the first
power of these ratios. Hence the right-hand side of
\eqref{eq:average} is at most
\[
 C_*(q^{-2/3}+q^{-1}+q^{-1/3}+q^{-1})
 \le4C_*q^{-1/3}=\frac12,
\]
where $q\ge1$ was used in the second inequality. This contradicts
\eqref{eq:average}.

Consequently no finite accepted prefix can have $K\ge s+m_*$:
if it did, its first $m_*$ accepted steps after $s$ would give the
contradiction just established. Prefixes with $K<s$ already satisfy
the same upper bound. Thus every finite accepted prefix obeys
\begin{equation}\label{eq:accepted-count}
 \begin{aligned}
 K<s+m_*=O\Biggl(&\Delta\left(
       \frac{\sqrt{L_0}}{\epsilon^{3/2}}
       +\frac{\sqrt{L_1}}{\epsilon}
       +\frac{L_0^2}{\delta^3}
       +\frac{L_1}{\delta}\right)\\
       &+1+\lp\frac{Q_0}{Q}
          +\lp\frac{\norm{\nabla f(x_0)}}{G_*}\Biggr).
 \end{aligned}
\end{equation}
The constants here are numerical and independent of all problem
parameters and tolerances. An infinite sequence of trials could
not contain only finitely many acceptances: after its last
acceptance, or from the initial state if no step were accepted,
the point would remain fixed, and Lemma~\ref{lem:safe}
would force acceptance after finitely many doublings unless the
algorithm had already stopped. An infinite run would therefore
have arbitrarily long finite accepted prefixes, contrary to
\eqref{eq:accepted-count}. The algorithm terminates in finitely
many trials, and its stopping test gives the required
$(\epsilon,\delta)$-second-order stationary point.

\paragraph{Rejected trials, oracle evaluations, and universality}
\label{sec:oracle-count}

If the initial point passes the stopping test, no trust-region
subproblem is solved and at most three initial oracle evaluations
are needed. Otherwise let $K\ge1$ and $J\ge0$ be the total numbers
of accepted and rejected trials before termination. The stopping
condition depends only on the current point, which does not change
on a rejection. Hence the algorithm stops immediately after its
last accepted step; no rejected trials follow that step.
Immediately before this last acceptance, the scale has been halved
exactly $K-1$ times and doubled exactly $J$ times. Thus
\[
 \sigma_{\tau_{K-1}}=\sigma_0\,2^{J-(K-1)},
 \qquad
 J=K-1+\log_2\frac{\sigma_{\tau_{K-1}}}{\sigma_0}.
\]
Using \eqref{eq:scale-global} and $Q_0\ge\sigma_0>0$, we obtain
\begin{align}
 J&\le K-1+\log_2C_E+\frac{1}{\log2}
                            \log\frac{Q_0}{\sigma_0}\nonumber\\
  &\le K+C\left(1+\lp\frac{Q_0}{\sigma_0}\right),
  \label{eq:rejected-count}
\end{align}
for a universal numerical constant $C$.

Every trial solves one trust-region subproblem, so the total
number of subproblems is $K+J$. The initial function value,
gradient, and Hessian cost at most three oracle evaluations.
Each trial evaluates one new function value and one new gradient;
on acceptance these are retained, and one new Hessian is evaluated
for the next stopping test or trial. On rejection the current data
are reused. Therefore the total number of exact function, gradient,
and Hessian evaluations, counting these three kinds together, is
at most
\[
 3+2(K+J)+K=3+3K+2J.
\]
Combining \eqref{eq:accepted-count}, \eqref{eq:rejected-count}, and
the definition of $\mathcal L$ in \eqref{eq:logs}, both counts are
bounded by
\begin{equation}\label{eq:explicit-total-count}
 C\left[\Delta\left(
       \frac{\sqrt{L_0}}{\epsilon^{3/2}}
       +\frac{\sqrt{L_1}}{\epsilon}
       +\frac{L_0^2}{\delta^3}
       +\frac{L_1}{\delta}\right)+\mathcal L\right].
\end{equation}
This also covers initial termination because $\mathcal L\ge1$.
It proves the stated bound, with the additive
$\widetilde O(1)$ term equal to $O(\mathcal L)$ and independent
of $\epsilon$ and $\delta$.

Finally, fix any pair $(L_0,L_1)$ satisfying the smoothness
assumption. Every preceding argument applies to that pair and
establishes \eqref{eq:explicit-total-count}. The algorithm uses
neither member of the pair, nor $f_*$ or $\Delta$; its trial
subproblems and acceptance decisions depend only on its inputs
and the oracle data. The same realized run thus satisfies the bound
for every admissible pair separately. This is the asserted
universality. All constants introduced for the potential and
counting arguments are proof constants and are not algorithmic
inputs.
